\section{Limitations and Broader Impact}
\label{sec:limitations}

\subsection{Limitations}

\paragraph{Scope.}
AmphionASR is optimised for Mandarin and English, with a small
amount of Mandarin--English code-switching data. We do not claim
broad multilingual coverage and we do not benchmark on languages
outside those reported in Section~\ref{sec:experiments}.

\paragraph{Open multi-talker benchmarks (positive WER).}
The target-speaker path is substantially weaker on the public
LibriMix test sets ($66.33$\,\% WER on Libri2Mix, $86.00$\,\% on
Libri3Mix) than on our internal \texttt{ts\_hw\_test} positive
sub-slice ($12.56$\,\% WER, Table~\ref{tab:tsasr}). The underlying
cause has not been investigated; until the next training
iteration the public multi-talker numbers should be read as a
known weakness, not a final result. The fact that Libri3Mix
degrades more than Libri2Mix relative to AmphionASR's own
\texttt{ts\_hw\_test} baseline (an extra $+19.67$ WER points
between two-speaker and three-speaker mixed audio) further
suggests that AmphionASR's speaker selection is not yet robust to
interferer count.

\paragraph{Negative-slice false alarms.}
On the $n_\text{neg} = 328$ target-absent samples of
\texttt{ts\_hw\_test}, AmphionASR reduces the false-alarm rate
from the open-baseline floor of $100\%$ to $98.17\%$
(Table~\ref{tab:tsasr}). The bulk of this reduction comes from
the \emph{silence} sub-type; on the \emph{distractor} sub-type
(only interfering speakers, no target) AmphionASR remains at
$100\%$ false alarms. The training-time augmentation source for
this stage -- Voices-in-the-Wild-2M -- contains no
\texttt{(audio, empty-target)} pairs, so the model receives no
explicit ``stay silent'' supervision; closing the gap requires
extending the SFT mixture with explicit anti-hallucination
negatives, which is the planned next step.

\paragraph{Noise robustness: residual gap on Voices-in-the-Wild-Bench.}
The Voices-in-the-Wild-Bench evaluation in
Section~\ref{sec:noise} reports AmphionASR as the best of nine
systems on five of sixteen subsets, with the remaining eleven
cells led by Mega-ASR~\cite{xie2026mega} (eight cells) or
Qwen3-ASR-1.7B~\cite{qwen2025qwen3asr} (three cells). The two
cells where AmphionASR trails Mega-ASR most clearly --
Sim-Echo \& reverberation and Sim-Recording coloration -- are
explicitly called out in Section~\ref{subsec:vitw-results}; the
underlying cause has not been investigated and closing the
residual is on the agenda for the next iteration. The eight
baseline columns in Table~\ref{tab:vitw_bench} are additionally
quoted under the Mega-ASR paper's normalisation pipeline rather
than ours, which is a known cross-system normalisation drift the
next revision is expected to retire by re-running each baseline.
% TODO: rerun all eight Voices-in-the-Wild-Bench baselines under
% our normaliser to retire the normalisation caveat in
% Section~\ref{sec:noise}.

\paragraph{Retrieve operator.}
The retrieval operator of
Section~\ref{subsec:retrieve-algorithm} has three known
boundary-of-applicability assumptions:
\begin{enumerate}
  \item The CJK pronunciation channel is keyed on tone-stripped
        Mandarin pinyin. Cantonese, Min, Hakka and the other
        Mandarin dialects covered by KeSpeech are not supported by
        this channel without replacing the pronunciation map; the
        framework is unchanged but the syllable inventory is
        Mandarin-only in this release.
  \item The alphabetic route assumes a content-word / stop-word
        decomposition that works for Indo-European languages. It
        has not been validated for agglutinative (Turkish,
        Finnish) or inflective languages, where the content-word
        overlap signal may degenerate.
  \item The text retrieval $r$ acts on the no-hotword coarse
        decode $\tilde{y}$. If $\tilde{y}$ contains a contiguous
        drop-out (long missing span) rather than the local
        substitution errors that the multi-signal score is
        designed for, $r$ can no longer recover the missing
        hotword and the operator degenerates to no-hotword ASR
        (the discriminative-gap check exists precisely to make
        this degeneration safe rather than catastrophic).
\end{enumerate}

\paragraph{Inherited LLM failure modes.}
As discussed in Section~\ref{sec:analysis}, the model can repeat on
long silences, drift between languages on code-switched content, and
over-bias on prompt-supplied hotwords when a term is not actually
spoken. These are well-known characteristics of LLM-based ASR and we
are explicit that AmphionASR shares them.

\paragraph{Latency and streaming.}
The release is offline-batch, served via vLLM. Streaming inference
and a measured first-token-latency report are out of scope for this
version.

\paragraph{Code release.}
We release the model weights and the vLLM-based serving stack that
produced every inference number in this report; the in-house
training code is not part of the release. The training pipeline is
documented end-to-end in Sections~\ref{sec:data}
and~\ref{sec:training} -- data sources, augmentation, three-stage
SFT and per-stage hyperparameters -- so that the recipe can be
reimplemented on top of a public training framework such as
ms-swift~\cite{msswift}. We adopt this weights-plus-serving release
scope because the in-house training code is tightly coupled with
internal data infrastructure and is not redistributable as it
stands; the scope is consistent with the practice of recent
SpeechLLM technical reports
\cite{radford2023whisper, kimiaudio2025, stepaudio22025,
qwen2025qwen3asr}.

\subsection{Broader Impact}

\paragraph{Surveillance risk.}
A target-speaker recognizer that operates from a short reference
utterance lowers the technical barrier to selectively transcribing a
designated person in multi-speaker audio. This capability has
legitimate uses (transcribing one's own voice in a meeting,
accessibility) but can also be misused for unconsented surveillance.
The released model card will state that the intended use is
user-initiated transcription with the consent of all recorded parties
and will explicitly call out unconsented surveillance as out-of-scope.

\paragraph{Bias.}
% TODO: insert per-accent / per-dialect WER breakdown when the
% KeSpeech-by-dialect split is added to the eval pack.
We have not yet measured per-accent WER systematically. KeSpeech is the
natural starting point for Mandarin dialectal coverage and will be
broken out in a follow-up.

\paragraph{Data provenance.}
% TODO: confirm provenance and licensing for every internal data source
% before public release of weights or checkpoints.
Public sources are used under their stated licenses (Tables~\ref{tab:data_asr}--\ref{tab:data_hotwords}).
Internal sources are subject to internal data-governance review prior
to any model-weight release.

\paragraph{LLM-assisted writing.}
Portions of this report were drafted with the assistance of large
language models. All technical claims, numbers, and citations were
verified by the authors against the corresponding primary sources.
% TODO: maintain ack/llm-usage.md as new sections are written and
% reflect the final disclosure here.
